\documentclass[11pt,a4paper]{article}

\usepackage[margin=2.0cm]{geometry}
\usepackage{booktabs}
\usepackage{multirow}
\usepackage{array}
\usepackage{xcolor}
\usepackage{hyperref}
\usepackage{parskip}
\usepackage{lmodern}
\usepackage{fancyvrb}

\hypersetup{colorlinks=true, linkcolor=black, urlcolor=blue}

\title{HyMeKo Generation \& Emission Framework\\[2pt]
       Multi-Target Codegen from a Single Hypergraph IR}
\author{thesis-IV working notes (HyMeKo)}
\date{2026-04-26}

\newcommand{\code}[1]{\texttt{\detokenize{#1}}}

\begin{document}
\maketitle

\begin{abstract}
HyMeKo describes robotic systems and neural-network architectures in
a single hypergraph IR. From the same description it can emit
URDF, SDF~1.7, MJCF, DOT, Mermaid, Gazebo world XML, and
torch\_dataflow PyTorch source. Each target is realised as a
\texttt{(queries.hymeko, template.X)} pair under
\texttt{transforms/<target>/}; the Rust template engine dispatches
queries against the IR and renders the template. We document the
architecture, the recently-added Tier~1 (plain MLP) and Tier~2
(residual / highway) emission paths for the PyTorch backend, and
the publication venues this systems-level contribution targets
separately from the entropy-regulariser paper.
\end{abstract}

\section{What can HyMeKo emit?}

Eight emission targets are currently shipped under
\code{transforms/}:

\begin{center}
\small
\begin{tabular}{@{}lll@{}}
\toprule
Target & Purpose & Status \\
\midrule
\code{urdf}            & ROS robot description       & shipped \\
\code{sdf}             & Gazebo scene description    & shipped \\
\code{mjcf}            & MuJoCo XML                  & shipped \\
\code{gazebo}          & Full Gazebo world XML       & shipped \\
\code{dot}             & Graphviz hypergraph viz     & shipped \\
\code{mermaid}         & Markdown-renderable graph   & shipped \\
\code{ros2\_launch}    & ROS~2 launch.py file        & shipped \\
\code{torch\_dataflow} & PyTorch \code{nn.Module}    & shipped + Tier 1 + Tier 2 (this report) \\
\bottomrule
\end{tabular}
\end{center}

A single \code{.hymeko} source can be compiled to any subset:

\begin{Verbatim}[fontsize=\small,frame=single]
hymeko compile data/robotics/anthropomorphic_arm.hymeko \
    --format urdf -o /tmp/arm.urdf

hymeko compile data/nn/mnist_small.hymeko \
    --format torch_dataflow --name MnistSmall -o /tmp/mnist_small.py
\end{Verbatim}

\section{Architecture: queries + template per target}

Every target is a directory \code{transforms/<target>/} with two
files:

\begin{description}[leftmargin=1.4em, style=nextline, itemsep=2pt]
  \item[\code{queries.hymeko}]  Declares named collections of decls
    matched by inheritance predicate. Example for
    \code{torch\_dataflow}:
    \begin{Verbatim}[fontsize=\small]
    layers:         meta_layer       {}
    linear_layers:  linear_layer     {}
    relu_layers:    relu_layer       {}
    residual_blocks: residual_block  {}
    @flows:         dataflow         {}
    outputs:        t_output         {}
    \end{Verbatim}
    Each line \texttt{name: type \{\}} becomes a collection that the
    template can iterate.

  \item[\code{template.<ext>}]  Render-time template using
    Handlebars-like syntax: \code{\{\{config:robot\_name\}\}} for
    config args, \code{\{\{\#each linear\_layers\}\}} to iterate
    matches, \code{\{\{name\}\}} / \code{\{\{field:d\_in\}\}} for
    per-decl values, \code{\{\{bind:+:0\}\}} for hyperedge arc
    bindings.
\end{description}

This separation has three consequences:

\begin{itemize}
  \item Adding a new emission target = adding a new directory.
        No changes to Rust code.
  \item Per-format expressivity is constrained to what the template
        engine + queries can express jointly. Multi-input dataflow
        and conditional emission per inheritance branch are
        currently supported; arithmetic on field values is not.
  \item The IR remains the single source of truth: structural
        consistency across emitted targets is automatic (a robot
        described in HyMeKo will have the same kinematic chain in its
        URDF and its MJCF and its DOT).
\end{itemize}

\section{Recent extension: Tier 1 + Tier 2 PyTorch dataflow}

Until 2026-04-26 the \code{torch\_dataflow} target only emitted
\code{HypergraphConv} layers --- HGNN-style networks. We extended it
to cover plain MLPs and residual / highway architectures so the
networks tested by the entropy-regulariser benchmark can also be
generated end-to-end from a HyMeKo description.

\subsection{Tier 1 — plain MLP primitives}

\textbf{New \code{meta\_layer} subtypes} added to
\code{data/nn/meta\_nn.hymeko}:

\begin{itemize}
  \item \code{linear\_layer}  --- emits \code{torch.nn.Linear(d\_in, d\_out)}
  \item \code{relu\_layer}    --- emits \code{torch.nn.ReLU()}
  \item \code{sigmoid\_layer} --- emits \code{torch.nn.Sigmoid()}
  \item \code{tanh\_layer}    --- emits \code{torch.nn.Tanh()}
\end{itemize}

\textbf{Example} (\code{data/nn/mnist\_small.hymeko}, 32 lines)
emits the same 784$\to$16$\to$8$\to$10 MLP that
\code{run\_benchmark.py}'s \code{MNISTNetSmall} hand-codes. Parameter
count: 12\,786 in both implementations (verified).

\subsection{Tier 2 — composite blocks}

For residual and gated networks, a fully-structural HyMeKo description
would need multi-input dataflow hyperedges (the additive skip needs
two incoming edges). That feature is deferred. Instead, Tier~2 adds
\textbf{named composite primitives} whose internal structure is
realised by helper classes shipped in
\code{python/ehk\_torch\_stub}:

\begin{itemize}
  \item \code{residual\_block}  --- emits \code{ResidualBlock(hidden=...)};
        helper class implements $y = x + \text{Linear}(\text{ReLU}(\text{Linear}(x)))$.
  \item \code{highway\_block}   --- emits \code{HighwayBlock(hidden=...)};
        helper implements $y = T(x) \odot F(x) + (1-T(x)) \odot x$.
\end{itemize}

\textbf{Example} (\code{data/nn/mnist\_resmlp\_3.hymeko}, 35 lines)
emits a 3-block ResMLP equivalent to
\code{ResMLP(hidden=16, n\_blocks=3)}. Parameter count: 14\,362 in
both implementations (verified).

\subsection{Why named-block instead of full-structural}

Three reasons:

\begin{enumerate}
  \item \textbf{Engineering scope.} Multi-input dataflow hyperedges
        are a real change to the template's iteration semantics
        (each \code{@flow} currently has a fixed (input, layer, output)
        shape). Adding multi-input is feasible but non-trivial.
  \item \textbf{Architectural clarity.} A named block is more
        readable than the equivalent hypergraph topology with explicit
        skip vertices. Domain experts describing networks at the
        \code{.hymeko} level want to write
        \code{block\_1: lyr.residual\_block { hidden 16; }}, not unwind
        the additive skip into hypergraph primitives.
  \item \textbf{Round-trippability.} Named blocks survive the
        round-trip from .hymeko $\to$ PyTorch $\to$ inspection. A
        flattened topology would lose the ``this is a residual''
        annotation needed for higher-level analysis.
\end{enumerate}

\section{Verification}

\begin{center}
\small
\begin{tabular}{@{}lrrl@{}}
\toprule
Description (lines) & Emitted module & Hand-coded equivalent & Param count match \\
\midrule
\code{mnist\_small.hymeko}    (32) & \code{MnistSmallFromHymeko} & \code{MNISTNetSmall} & 12\,786 / 12\,786 \,$\checkmark$ \\
\code{mnist\_resmlp\_3.hymeko} (35) & \code{ResMLP3FromHymeko}    & \code{ResMLP(16, 3)} & 14\,362 / 14\,362 \,$\checkmark$ \\
\bottomrule
\end{tabular}
\end{center}

The emitted networks instantiate cleanly and forward-pass end-to-end
on a $4 \times 784$ input batch, returning the expected
$4 \times 10$ output.

\section{What this is not yet}

\begin{itemize}
  \item \textbf{No CapsMLP emission.}  Dynamic routing iterations
        require either inline Python loop emission or a fully
        opaque \code{CapsuleRouting} helper. Both are doable; deferred.
  \item \textbf{No batch-norm / dropout emission.}  These need to be
        decided at the \code{linear\_layer} level (annotation field?)
        or as separate single-purpose layer types
        (\code{batchnorm\_layer}, \code{dropout\_layer}). Trivial to
        add when the use case requires.
  \item \textbf{No multi-input dataflow hyperedges.}  Required for
        fully-structural ResNet / Highway / U-Net descriptions.
        Targeted for v0.3.
  \item \textbf{No regulariser-aware emission.}  The emitted PyTorch
        modules don't expose the \code{spectral\_weights()} method
        the entropy regulariser needs. Could be auto-emitted from the
        \code{linear\_layer} list. Trivial; deferred.
\end{itemize}

\section{Publication venues for the emission framework}

Distinct from the entropy-regulariser paper (which is a methods/
optimisation contribution at TMLR/NeurIPS/AISTATS scope). The
generation/emission framework is a \emph{language and systems}
contribution. Realistic venues:

\begin{description}[leftmargin=1.4em, style=nextline, itemsep=2pt]
  \item[MDPI Technologies (current).] Already submitting the
        hypergraph IR for robotic systems. Tier 1 / Tier 2 PyTorch
        emission can be added as a methods extension or a follow-up
        application paper.
  \item[Software Impacts (Elsevier).] Short-form journal explicitly
        for computer-science software contributions. The multi-target
        codegen architecture fits the format exactly.
  \item[JOSS (Journal of Open Source Software).] If we publish a
        version of HyMeKo as a standalone open-source release. Short
        paper format ($\approx$\,1000 words), the emphasis is on the
        software's utility and reproducibility.
  \item[MODELS / ACM SIGPLAN DSL workshops.] Domain-specific
        languages and model-driven engineering. The HyMeKo grammar +
        emission targets fit this audience well.
  \item[ICRA / IROS / RSS — workshop tracks.] For the robotic-
        systems angle (URDF / SDF / MJCF emission). The Galambos
        student work on VLM-driven robot descriptions could feed an
        ICRA workshop submission.
  \item[ICLR / NeurIPS workshops on compositional AI / DSL for ML.]
        The PyTorch emission story (especially with regulariser-aware
        emission added) is a reasonable workshop fit.
\end{description}

\textbf{Strongest recommendation:} Software Impacts plus a JOSS paper
in parallel. Both are short, both reward reproducibility (already met
by the workspace's reproducible-test infrastructure), and both
provide citable artefacts that downstream ML/robotics work can refer
to. The MDPI Technologies and ICRA-workshop venues are application
papers built on top of this systems contribution.

\section{Architecture summary}

The full emission stack (single direction):

\begin{Verbatim}[fontsize=\small,frame=single]
.hymeko source
    |
    | parser (LALRPOP, lexer with SIMD)
    v
AST (parser::ast::AstStr<'a>)
    |
    | hymeko_core::module_store::ModuleStore::compile
    v
Compiled IR (CompiledProgram: nodes, edges, arcs, decl tree, hash)
    |
    | hymeko_query::transforms::TransformRegistry
    | + transforms/<target>/queries.hymeko
    v
Per-format query results (Vec<QueryMatch>)
    |
    | template engine
    | + transforms/<target>/template.<ext>
    v
Emitted artefact (URDF, MJCF, .py, .dot, ...)
\end{Verbatim}

The hot extension surface (where Tier 1 / Tier 2 changes landed) is
on the right edge: the \code{queries.hymeko} and \code{template.py}
files for the \code{torch\_dataflow} format. Zero changes to Rust
code; the extension is purely declarative.

\section{Reproducibility}

To regenerate both example emissions from this workspace:

\begin{Verbatim}[fontsize=\small,frame=single]
target/release/hymeko compile data/nn/mnist_small.hymeko \
    --format torch_dataflow --name MnistSmallFromHymeko \
    -o /tmp/mnist_small_emitted.py

target/release/hymeko compile data/nn/mnist_resmlp_3.hymeko \
    --format torch_dataflow --name ResMLP3FromHymeko \
    -o /tmp/mnist_resmlp_emitted.py

PYTHONPATH=python/ehk_torch_stub/src python3 -c "
import torch, importlib.util, sys
def load(p, n):
    s = importlib.util.spec_from_file_location(n, p)
    m = importlib.util.module_from_spec(s); s.loader.exec_module(m); return m
m1 = load('/tmp/mnist_small_emitted.py', 'm1').MnistSmallFromHymeko()
m2 = load('/tmp/mnist_resmlp_emitted.py', 'm2').ResMLP3FromHymeko()
print('p1 =', sum(p.numel() for p in m1.parameters()),
      ', p2 =', sum(p.numel() for p in m2.parameters()))
"
# Expected:  p1 = 12786 , p2 = 14362
\end{Verbatim}

\end{document}
